% Generated draft: requires mathematical and source review.
\documentclass[11pt]{article}
\usepackage[margin=25mm]{geometry}
\usepackage{fontspec}
\setmainfont{Latin Modern Roman}
\usepackage{amsmath,amssymb,mathtools,unicode-math}
\setmathfont{Latin Modern Math}
\usepackage{enumitem,needspace,xurl,hyperref}
\setlength{\parindent}{0pt}
\setlength{\parskip}{5pt}
\setlength{\emergencystretch}{3em}
\newcommand{\N}{\mathbb{N}}
\newcommand{\Z}{\mathbb{Z}}
\newcommand{\Q}{\mathbb{Q}}
\newcommand{\R}{\mathbb{R}}
\newcommand{\C}{\mathbb{C}}
\newcommand{\F}{\mathbb{F}}
\newcommand{\A}{\mathbb{A}}
\newcommand{\PP}{\mathbb{P}}
\newcommand{\E}{\mathbb{E}}
\newcommand{\eps}{\varepsilon}
\newcommand{\dd}{\,\mathrm{d}}
\begin{document}
\section*{UnsolvedMath ID: 30001195}
\textbf{Source problem code:} catalogue entry
\subsection{Definitions and mathematical statement}
The supplied question concerns a Lie superalgebra \(\mathfrak g=\mathfrak g_{\bar 0}\oplus\mathfrak g_{\bar 1}\), a simple highest-weight \(\mathfrak g\)-module \(L(\lambda)\), and an associated Lie subsuperalgebra \(\mathfrak e\subseteq\mathfrak g\), whose even part is \(\mathfrak e_{\bar 0}\). The record does not specify the ground field, the precise class of Lie superalgebras, the module category, or the construction of \(\mathfrak e\). Consequently, the following definitions describe a conventional interpretation conditional on those choices, rather than supplying missing source hypotheses.

In a root-theoretic setting where the usual definition of atypicality applies, fix positive even and odd roots \(\Delta_{\bar 0}^{+}\) and \(\Delta_{\bar 1}^{+}\), together with the bilinear form \((\cdot,\cdot)\) on weights used in that theory, and put
\[
\rho=\frac12\sum_{\alpha\in\Delta_{\bar 0}^{+}}\alpha-\frac12\sum_{\beta\in\Delta_{\bar 1}^{+}}\beta.
\]
The atypicality \(\operatorname{atyp}(\lambda)\) is the largest integer \(r\geq 0\) for which there are linearly independent odd roots \(\beta_1,\ldots,\beta_r\) satisfying
\[
(\beta_i,\beta_j)=0\quad(1\leq i,j\leq r),
\qquad
(\lambda+\rho,\beta_i)=0\quad(1\leq i\leq r).
\]
In particular, each \(\beta_i\) is isotropic. The root data and hypotheses needed for this interpretation must be confirmed from the cited source.

For a relative cohomological interpretation, write
\[
R_{\mathfrak e}=H^{\bullet}(\mathfrak e,\mathfrak e_{\bar 0};k),
\qquad
I_M=\operatorname{Ann}_{R_{\mathfrak e}}
\operatorname{Ext}^{\bullet}_{(\mathfrak e,\mathfrak e_{\bar 0})}(M,M),
\]
where \(k\) is the ground field and the action is the usual relative cohomology action. Whenever the intended theory supplies this commutative, finitely generated coordinate ring, its cohomological support variety is
\[
V_{(\mathfrak e,\mathfrak e_{\bar 0})}(M)
=\operatorname{MaxSpec}(R_{\mathfrak e}/I_M),
\]
with dimension understood as Krull dimension. For \(M=L(\lambda)\), this construction uses restriction to \(\mathfrak e\). The supplied wording also calls the variety a rank variety; identifying these constructions requires hypotheses and an identification not given in the record.

The source question is whether, for every simple highest-weight module in the intended scope,
\[
\operatorname{atyp}(\lambda)
=\dim V_{(\mathfrak e,\mathfrak e_{\bar 0})}\bigl(L(\lambda)\bigr).
\]
The universal quantifier is over the source's intended Lie superalgebras and modules, whose precise scope remains unresolved. The supplied literature assessment mentions both \(\mathfrak{gl}(m\mid n)\) and simple basic classical Lie superalgebras, but does not provide the hypotheses needed to turn these mentions into a fully specified common formulation.
\subsection{Short English statement}
The question asks whether the atypicality of each intended simple highest-weight supermodule equals the dimension of its associated support or rank variety. The ground field, module category, class of Lie superalgebras, and choice of \(\mathfrak e\) require clarification. Supplied dated curation labels the problem partially solved, reporting the equality for \(\mathfrak{gl}(m\mid n)\) and a conjectural or case-dependent situation for all simple basic classical Lie superalgebras; these claims are not independently verified here.
\subsection{Sources}
\begin{itemize}
\item Supplied local UnsolvedMath registry record, permanent ID \(30001195\), title: Atypicality and Support Varieties of Simple Supermodules. Exact source problem code: \texttt{catalogue entry}. Collection: Oberwolfach Reports Open Problems; proposed year: 2009. No individual proposer is identified.
\item Supplied original report citation: Mini-Workshop: Support Varieties (2009), Oberwolfach Reports, pp.\ 549--592, DOI \texttt{10.4171/owr/2009/10}. \url{https://doi.org/10.4171/owr/2009/10}. The record describes these reports as documenting mathematical workshops and their open-problem discussions.
\item Supplied literature reference labeled Boe--Kujawa--Nakano II: \url{https://arxiv.org/abs/0708.3191}.
\item Supplied literature reference labeled published primary result: \url{https://doi.org/10.1112/plms/pdn019}.
\item Supplied literature reference labeled original Oberwolfach problem: \url{https://doi.org/10.4171/owr/2009/10}.
\item Literature assessment date: \texttt{2026-08-21}. This is supplied dated curation derived from the supplied OWR curation, not an independently verified literature assessment. It labels the status partially solved and reports that atypicality equals support dimension for \(\mathfrak{gl}(m\mid n)\), while the corresponding assertion for all simple basic classical Lie superalgebras remains conjectural or case-dependent. The record requests independent verification of time-sensitive status and cited claims.
\item The registry labels the statement \texttt{corrected\_verified} and reports that its mathematical content was checked against the official Oberwolfach report, with notation, context, and question form normalized without changing the source-level claim. This is supplied verification metadata; no source visit or independent verification is asserted here.
\end{itemize}
\end{document}
